% Purpose: the screening demonstration (claim C5): out-of-sample predicted against DFT neutral oxygen-vacancy formation energies of named oxides, both backbones, with the scope measured and the chemistry discussed. Every number is a macro of paper/generated/numbers.tex (sources: figures/main/data/screening_demo.json, results/screening_absence.json).
\subsection{Screening demonstration}
\label{sec:screening}

The networks were built to predict, from the crystal structure alone, the energy needed to remove one neutral oxygen atom from an oxide. This subsection applies the structure-only network S to named oxides whose DFT formation energies are in the release, so that every prediction can be set against the value of the same site. It is a demonstration, not a test of a hypothesis (claim C5). Each oxide shown was a test host in at least one outer resample, so none of its sites was in the training set behind its prediction, and the budget is \numScreeningDemoBudget{} training hosts, the largest of the learning curves. The prediction for a site is the seed-ensemble mean averaged over the outer resamples in which its host was held out, the aggregation fixed in the analysis plan; Figure~\ref{fig:screening} shows, as a bar, the smallest and the largest of the per-resample predictions. The candidate list and this aggregation were fixed before any prediction was made, and the originally registered backbone (v1) is the backbone the plan describes. The primary backbone (v2) was adopted after diagnostics of v1, so its predictions are post hoc. Both are shown.

\subsubsection{Which oxides are present, and which are not.}
The plan names \numScreeningDemoScopeNCandidates{} candidate oxides. \numScreeningDemoScopeNPresent{} are in universe v1: ZnO, TiO$_2$, SnO$_2$, MgO, Al$_2$O$_3$, ZrO$_2$ and Ga$_2$O$_3$, with \numScreeningDemoScopeNScreenedSites{} oxygen sites in \numScreeningDemoScopeNHosts{} hosts (Table~\ref{tab:screening_sites}). The other \numScreeningDemoScopeNAbsent{} are absent, and the release shows two different reasons (\texttt{results/screening\_absence.json}). \numScreeningDemoScopeNNotInRelease{} of them, In$_2$O$_3$, CeO$_2$, WO$_3$, MoO$_3$, Nb$_2$O$_5$ and LaAlO$_3$, have no host of that formula in the release; the release does not record why, and we state no reason. The remaining \numScreeningDemoScopeNPhsRemoved{}, SrTiO$_3$ and BaTiO$_3$, are in the release, but every one of their neutral entries carries the perturbed-host-state flag (\numScreeningDemoScopePhsRemovedSrTiOThreePhsEntries{} of \numScreeningDemoScopePhsRemovedSrTiOThreeNeutralEntries{} for SrTiO$_3$, \numScreeningDemoScopePhsRemovedBaTiOThreePhsEntries{} of \numScreeningDemoScopePhsRemovedBaTiOThreeNeutralEntries{} for BaTiO$_3$), so the filter of Section~\ref{sec:data} leaves no site of either. TiO$_2$ stays with a single site, because \numScreeningDemoScopeTioTwoReleasePhsEntries{} of its \numScreeningDemoScopeTioTwoReleaseNeutralEntries{} neutral entries are flagged. The release host of TiO$_2$ is a monoclinic polymorph (space group number \numScreeningDemoHostsTiOTwoSpaceGroupNumber{}, from the supercell of the release), not rutile. That the flag removes the titanates and most of the TiO$_2$ sites is consistent with the hypothesis of Section~\ref{sec:data} that oxides with a low-lying empty cation $d$ band are over-represented among shallow-donor entries; it does not prove it. The demonstration therefore covers only oxides whose neutral vacancy the release treats as a localised defect, and it says nothing about the absent ones. Of the \numScreeningDemoScopeNScreenedSites{} sites, \numScreeningDemoScopeNSitesOneResample{} were held out in a single resample, and their predictions have no spread.

\begin{table}[H]
\caption{Screened oxygen sites. Structure: textbook structure type, assigned from the space group and the coordination of the site, or monoclinic where no simple type applies; SG: space group number of the host supercell. Cations: number of cations in the first shell of the vacancy oxygen, the cations within a fixed distance factor of the nearest one (the shell of the structure renders, \texttt{blender/README.md}); $\langle d\rangle$: their mean distance to the oxygen. $E_\mathrm{g}$: host band gap and $\Delta H_\mathrm{f}$: host formation energy per atom, both from the descriptors of the release. DFT: formation energy of the release. Predictions of S: primary backbone (v2) and originally registered backbone (v1), mean over the $N$ outer resamples in which the host was held out. Source: \texttt{figures/main/data/screening\_demo.json}.\label{tab:screening_sites}}
\begin{adjustwidth}{-\extralength}{0cm}
\centering
\small
\begin{tabular}{llllrrrrrrrr}
\toprule
Oxide & Site & Structure & SG & Cations & $\langle d\rangle$ (\AA) & $E_\mathrm{g}$ (eV) & $\Delta H_\mathrm{f}$ (eV) & DFT (eV) & v2 (eV) & v1 (eV) & $N$ \\
\midrule
SnO$_2$ & O$_1$ & rutile-type & \numScreeningDemoHostsSnOTwoSpaceGroupNumber & \numScreeningDemoSitesSnOTwoVaOOneZeroFirstShellCations & \numScreeningDemoSitesSnOTwoVaOOneZeroFirstShellMeanDistanceA & \numScreeningDemoSitesSnOTwoVaOOneZeroBandGapEV & \numScreeningDemoSitesSnOTwoVaOOneZeroHostFormationEnergyEVPerAtom & \numScreeningDemoSitesSnOTwoVaOOneZeroDftEV & \numScreeningDemoSitesSnOTwoVaOOneZeroVTwoPredEV & \numScreeningDemoSitesSnOTwoVaOOneZeroVOnePredEV & \numScreeningDemoSitesSnOTwoVaOOneZeroVTwoNResamples \\
ZnO & O$_1$ & wurtzite & \numScreeningDemoHostsZnOSpaceGroupNumber & \numScreeningDemoSitesZnOVaOOneZeroFirstShellCations & \numScreeningDemoSitesZnOVaOOneZeroFirstShellMeanDistanceA & \numScreeningDemoSitesZnOVaOOneZeroBandGapEV & \numScreeningDemoSitesZnOVaOOneZeroHostFormationEnergyEVPerAtom & \numScreeningDemoSitesZnOVaOOneZeroDftEV & \numScreeningDemoSitesZnOVaOOneZeroVTwoPredEV & \numScreeningDemoSitesZnOVaOOneZeroVOnePredEV & \numScreeningDemoSitesZnOVaOOneZeroVTwoNResamples \\
Ga$_2$O$_3$ & O$_1$ & monoclinic & \numScreeningDemoHostsGaTwoOThreeSpaceGroupNumber & \numScreeningDemoSitesGaTwoOThreeVaOOneZeroFirstShellCations & \numScreeningDemoSitesGaTwoOThreeVaOOneZeroFirstShellMeanDistanceA & \numScreeningDemoSitesGaTwoOThreeVaOOneZeroBandGapEV & \numScreeningDemoSitesGaTwoOThreeVaOOneZeroHostFormationEnergyEVPerAtom & \numScreeningDemoSitesGaTwoOThreeVaOOneZeroDftEV & \numScreeningDemoSitesGaTwoOThreeVaOOneZeroVTwoPredEV & \numScreeningDemoSitesGaTwoOThreeVaOOneZeroVOnePredEV & \numScreeningDemoSitesGaTwoOThreeVaOOneZeroVTwoNResamples \\
Ga$_2$O$_3$ & O$_2$ & monoclinic & \numScreeningDemoHostsGaTwoOThreeSpaceGroupNumber & \numScreeningDemoSitesGaTwoOThreeVaOTwoZeroFirstShellCations & \numScreeningDemoSitesGaTwoOThreeVaOTwoZeroFirstShellMeanDistanceA & \numScreeningDemoSitesGaTwoOThreeVaOTwoZeroBandGapEV & \numScreeningDemoSitesGaTwoOThreeVaOTwoZeroHostFormationEnergyEVPerAtom & \numScreeningDemoSitesGaTwoOThreeVaOTwoZeroDftEV & \numScreeningDemoSitesGaTwoOThreeVaOTwoZeroVTwoPredEV & \numScreeningDemoSitesGaTwoOThreeVaOTwoZeroVOnePredEV & \numScreeningDemoSitesGaTwoOThreeVaOTwoZeroVTwoNResamples \\
Ga$_2$O$_3$ & O$_3$ & monoclinic & \numScreeningDemoHostsGaTwoOThreeSpaceGroupNumber & \numScreeningDemoSitesGaTwoOThreeVaOThreeZeroFirstShellCations & \numScreeningDemoSitesGaTwoOThreeVaOThreeZeroFirstShellMeanDistanceA & \numScreeningDemoSitesGaTwoOThreeVaOThreeZeroBandGapEV & \numScreeningDemoSitesGaTwoOThreeVaOThreeZeroHostFormationEnergyEVPerAtom & \numScreeningDemoSitesGaTwoOThreeVaOThreeZeroDftEV & \numScreeningDemoSitesGaTwoOThreeVaOThreeZeroVTwoPredEV & \numScreeningDemoSitesGaTwoOThreeVaOThreeZeroVOnePredEV & \numScreeningDemoSitesGaTwoOThreeVaOThreeZeroVTwoNResamples \\
TiO$_2$ & O$_4$ & monoclinic & \numScreeningDemoHostsTiOTwoSpaceGroupNumber & \numScreeningDemoSitesTiOTwoVaOFourZeroFirstShellCations & \numScreeningDemoSitesTiOTwoVaOFourZeroFirstShellMeanDistanceA & \numScreeningDemoSitesTiOTwoVaOFourZeroBandGapEV & \numScreeningDemoSitesTiOTwoVaOFourZeroHostFormationEnergyEVPerAtom & \numScreeningDemoSitesTiOTwoVaOFourZeroDftEV & \numScreeningDemoSitesTiOTwoVaOFourZeroVTwoPredEV & \numScreeningDemoSitesTiOTwoVaOFourZeroVOnePredEV & \numScreeningDemoSitesTiOTwoVaOFourZeroVTwoNResamples \\
ZrO$_2$ & O$_1$ & monoclinic & \numScreeningDemoHostsZrOTwoSpaceGroupNumber & \numScreeningDemoSitesZrOTwoVaOOneZeroFirstShellCations & \numScreeningDemoSitesZrOTwoVaOOneZeroFirstShellMeanDistanceA & \numScreeningDemoSitesZrOTwoVaOOneZeroBandGapEV & \numScreeningDemoSitesZrOTwoVaOOneZeroHostFormationEnergyEVPerAtom & \numScreeningDemoSitesZrOTwoVaOOneZeroDftEV & \numScreeningDemoSitesZrOTwoVaOOneZeroVTwoPredEV & \numScreeningDemoSitesZrOTwoVaOOneZeroVOnePredEV & \numScreeningDemoSitesZrOTwoVaOOneZeroVTwoNResamples \\
ZrO$_2$ & O$_2$ & monoclinic & \numScreeningDemoHostsZrOTwoSpaceGroupNumber & \numScreeningDemoSitesZrOTwoVaOTwoZeroFirstShellCations & \numScreeningDemoSitesZrOTwoVaOTwoZeroFirstShellMeanDistanceA & \numScreeningDemoSitesZrOTwoVaOTwoZeroBandGapEV & \numScreeningDemoSitesZrOTwoVaOTwoZeroHostFormationEnergyEVPerAtom & \numScreeningDemoSitesZrOTwoVaOTwoZeroDftEV & \numScreeningDemoSitesZrOTwoVaOTwoZeroVTwoPredEV & \numScreeningDemoSitesZrOTwoVaOTwoZeroVOnePredEV & \numScreeningDemoSitesZrOTwoVaOTwoZeroVTwoNResamples \\
MgO & O$_1$ & rocksalt & \numScreeningDemoHostsMgOSpaceGroupNumber & \numScreeningDemoSitesMgOVaOOneZeroFirstShellCations & \numScreeningDemoSitesMgOVaOOneZeroFirstShellMeanDistanceA & \numScreeningDemoSitesMgOVaOOneZeroBandGapEV & \numScreeningDemoSitesMgOVaOOneZeroHostFormationEnergyEVPerAtom & \numScreeningDemoSitesMgOVaOOneZeroDftEV & \numScreeningDemoSitesMgOVaOOneZeroVTwoPredEV & \numScreeningDemoSitesMgOVaOOneZeroVOnePredEV & \numScreeningDemoSitesMgOVaOOneZeroVTwoNResamples \\
Al$_2$O$_3$ & O$_1$ & corundum & \numScreeningDemoHostsAlTwoOThreeSpaceGroupNumber & \numScreeningDemoSitesAlTwoOThreeVaOOneZeroFirstShellCations & \numScreeningDemoSitesAlTwoOThreeVaOOneZeroFirstShellMeanDistanceA & \numScreeningDemoSitesAlTwoOThreeVaOOneZeroBandGapEV & \numScreeningDemoSitesAlTwoOThreeVaOOneZeroHostFormationEnergyEVPerAtom & \numScreeningDemoSitesAlTwoOThreeVaOOneZeroDftEV & \numScreeningDemoSitesAlTwoOThreeVaOOneZeroVTwoPredEV & \numScreeningDemoSitesAlTwoOThreeVaOOneZeroVOnePredEV & \numScreeningDemoSitesAlTwoOThreeVaOOneZeroVTwoNResamples \\
\bottomrule
\end{tabular}
\end{adjustwidth}
\end{table}

\begin{figure}[H]
\figph{fig7_screening}
\caption{Out-of-sample predicted and DFT neutral oxygen-vacancy formation energies of the screened oxides, one marker per site and colour per oxide. (a)~Predicted against DFT with the line of equality. (b)~Predicted minus DFT, grouped by oxide in ascending order of its mean DFT energy. Filled markers and solid bars: primary backbone (v2); open markers and dashed bars: originally registered backbone (v1). A bar spans the smallest to the largest per-resample prediction of the outer resamples in which the host was held out; a site held out once has no bar. Sources: \texttt{figures/main/data/screening\_demo.json} and \texttt{results/screening\_c5.json}.\label{fig:screening}}
\end{figure}

\subsubsection{Screening error.}
Over the \numScreeningDemoScopeNScreenedSites{} sites the mean absolute error is \numScreeningDemoSummaryVOneMaeEV{}~eV for v1 (two-sided interval at level \numCfgStatsCi{}, \numScreeningDemoSummaryVOneMaeCiLo{} to \numScreeningDemoSummaryVOneMaeCiHi{}~eV) and \numScreeningDemoSummaryVTwoMaeEV{}~eV for v2 (\numScreeningDemoSummaryVTwoMaeCiLo{} to \numScreeningDemoSummaryVTwoMaeCiHi{}~eV). The intervals are percentile bootstrap intervals over hosts, the sites of one oxide forming one cluster; with \numScreeningDemoScopeNHosts{} hosts they are wide, and they are not the hierarchical bootstrap of the learning curves. The result file of the registered screening (\texttt{results/screening\_c5.json}) holds the two point values and no interval. For comparison, the error of the same networks over all held-out sites at the same budget is \numCOneCTwoVOneMaeSixHundredFiftyFourSMaeMean{}~eV for v1 (\numCOneCTwoVOneMaeSixHundredFiftyFourSMaeCiLo{} to \numCOneCTwoVOneMaeSixHundredFiftyFourSMaeCiHi{}~eV) and \numCOneCTwoVTwoMaeSixHundredFiftyFourSMaeMean{}~eV for v2 (\numCOneCTwoVTwoMaeSixHundredFiftyFourSMaeCiLo{} to \numCOneCTwoVTwoMaeSixHundredFiftyFourSMaeCiHi{}~eV). The v1 screening error lies inside its all-site interval; the v2 screening error lies above its all-site interval, although the two intervals overlap. The screened oxides are not a random sample of hosts, and we did not test whether the screening error differs from the all-site error. On these sites v1 has the smaller error; v2 is the primary backbone for the whole sweep, and we do not read a ranking of backbones from ten sites. The errors have no net sign: for each backbone \numScreeningDemoSummaryVTwoNOverPrediction{} sites are over-predicted and \numScreeningDemoSummaryVTwoNUnderPrediction{} under-predicted, and the mean signed error is \numScreeningDemoSummaryVTwoBiasEV{}~eV for v2 and \numScreeningDemoSummaryVOneBiasEV{}~eV for v1.

\subsubsection{The energy scale across oxides.}
Panel~(a) of Figure~\ref{fig:screening} shows two groups of DFT energies: the sites of SnO$_2$, ZnO, Ga$_2$O$_3$ and TiO$_2$ lie lower, those of ZrO$_2$, MgO and Al$_2$O$_3$ higher, and both backbones place every site in its group. Across the \numScreeningDemoScopeNScreenedSites{} sites the Spearman rank correlation of the DFT energy with the predicted energy is \numScreeningDemoSummaryVTwoSpearmanAcrossSites{} for v2 and \numScreeningDemoSummaryVOneSpearmanAcrossSites{} for v1; this is dominated by the separation of the two groups. Removing a neutral oxygen atom breaks the cation--oxygen bonds around the site and leaves its two electrons on the neighbouring cations, so the cost should grow with the strength of the host's bonding and with the energy at which the electrons can be accommodated. The data are descriptively in line with this: the DFT energy of the screened sites is rank-correlated with the host band gap (Spearman correlation \numScreeningDemoDescriptiveSpearmanDftVsBandGap{}) and anticorrelated with the host formation energy per atom (\numScreeningDemoDescriptiveSpearmanDftVsHostFormationEnergy{}, so a more negative formation energy goes with a higher vacancy energy). These are ten sites in seven hosts, the sites of one host are not independent, and no interval is attached. Hypothesis: in the wide-gap, strongly bound oxides MgO, Al$_2$O$_3$ and ZrO$_2$ the bonds are costly to break and the electrons left behind must occupy a defect level high in a wide gap; in ZnO, SnO$_2$ and Ga$_2$O$_3$, whose conduction bands are of cation $s$ character and whose gaps are smaller (Table~\ref{tab:screening_sites}), they can be placed at lower cost. TiO$_2$ falls in the lower group although its host formation energy per atom is more negative than that of MgO (Table~\ref{tab:screening_sites}), which bond strength alone does not explain. It is also the only screened oxide with flagged sites in the release, which places it nearest the edge of the domain. Neither statement was tested here.

\subsubsection{The order of sites within one oxide.}
Two hosts have more than one screened site (Table~\ref{tab:screening_pairs}). In Ga$_2$O$_3$ the DFT energy is highest on the site with four neighbouring cations (O$_2$) and lower on the two sites with three (O$_1$ and O$_3$), which is the order that counting the broken bonds gives. Both backbones order all three Ga$_2$O$_3$ pairs as DFT does. The separation is not resolved, however: for the pair O$_1$ and O$_3$, which bond counting cannot separate, the DFT energies differ by the amount in Table~\ref{tab:screening_pairs} and the predicted difference is a small fraction of it, and the per-resample ranges of the two sites overlap in both backbones (Figure~\ref{fig:screening}b). The network separates the four-fold site from the three-fold sites, but in this sample it does not separate the two three-fold sites from each other. In ZrO$_2$ the order is wrong in both backbones: DFT places the four-fold site O$_2$ slightly below the three-fold site O$_1$, against bond counting, and the predictions place it at or above, as bond counting would. The mean cation--oxygen distance of O$_2$ is the longer of the two (Table~\ref{tab:screening_sites}). Hypothesis: the longer, weaker bonds outweigh the larger number of bonds in ZrO$_2$ but not in Ga$_2$O$_3$, where the four-fold site also has the longer bonds and the highest energy, so that neither the coordination nor the mean distance alone orders the sites of both oxides. Within the \numScreeningDemoSummaryVTwoWithinHostPairs{} within-host pairs both backbones get \numScreeningDemoSummaryVTwoWithinHostPairsConcordant{} in the DFT order.

\begin{table}[H]
\caption{Differences between two screened sites of one host, first site minus second, from DFT and from the predictions of S (primary backbone v2 and originally registered backbone v1; mean over the outer resamples in which the host was held out). A prediction has the order of DFT when the difference has the sign of the DFT difference. Source: \texttt{figures/main/data/screening\_demo.json}, keys \texttt{pairs}.\label{tab:screening_pairs}}
\centering
\small
\begin{tabular}{llrrr}
\toprule
Oxide & Sites & DFT (eV) & v2 (eV) & v1 (eV) \\
\midrule
Ga$_2$O$_3$ & O$_1$$-$O$_2$ & \numScreeningDemoPairsGaTwoOThreeOOneOTwoDftDiffEV & \numScreeningDemoPairsGaTwoOThreeOOneOTwoVTwoDiffEV & \numScreeningDemoPairsGaTwoOThreeOOneOTwoVOneDiffEV \\
Ga$_2$O$_3$ & O$_1$$-$O$_3$ & \numScreeningDemoPairsGaTwoOThreeOOneOThreeDftDiffEV & \numScreeningDemoPairsGaTwoOThreeOOneOThreeVTwoDiffEV & \numScreeningDemoPairsGaTwoOThreeOOneOThreeVOneDiffEV \\
Ga$_2$O$_3$ & O$_2$$-$O$_3$ & \numScreeningDemoPairsGaTwoOThreeOTwoOThreeDftDiffEV & \numScreeningDemoPairsGaTwoOThreeOTwoOThreeVTwoDiffEV & \numScreeningDemoPairsGaTwoOThreeOTwoOThreeVOneDiffEV \\
ZrO$_2$ & O$_1$$-$O$_2$ & \numScreeningDemoPairsZrOTwoOOneOTwoDftDiffEV & \numScreeningDemoPairsZrOTwoOOneOTwoVTwoDiffEV & \numScreeningDemoPairsZrOTwoOOneOTwoVOneDiffEV \\
\bottomrule
\end{tabular}
\end{table}

\subsubsection{Where the model is wrong.}
The largest error of v2 is \numScreeningDemoSummaryVTwoMaxAbsErrEV{}~eV, at the O$_2$ site of ZrO$_2$; that of v1 is \numScreeningDemoSummaryVOneMaxAbsErrEV{}~eV, at the O$_4$ site of TiO$_2$. MgO and Al$_2$O$_3$, the two sites of highest DFT energy, are under-predicted by both backbones, v2 by \numScreeningDemoSitesMgOVaOOneZeroVTwoAbsErrEV{} and \numScreeningDemoSitesAlTwoOThreeVaOOneZeroVTwoAbsErrEV{}~eV (MgO and Al$_2$O$_3$) and v1 by \numScreeningDemoSitesMgOVaOOneZeroVOneAbsErrEV{} and \numScreeningDemoSitesAlTwoOThreeVaOOneZeroVOneAbsErrEV{}~eV; both ZrO$_2$ sites and the TiO$_2$ site are over-predicted by both. The sign of the error therefore follows no single rule across the oxides. The spread over resamples is a property of the site and of the backbone, and neither backbone has the smaller spread everywhere. For the O$_2$ site of Ga$_2$O$_3$ the predictions of v1 run from \numScreeningDemoSitesGaTwoOThreeVaOTwoZeroVOneMinEV{} to \numScreeningDemoSitesGaTwoOThreeVaOTwoZeroVOneMaxEV{}~eV across resamples against a DFT value of \numScreeningDemoSitesGaTwoOThreeVaOTwoZeroDftEV{}~eV, so the small error of the mean (Table~\ref{tab:screening_sites}) hides single-resample errors of several tenths of an electronvolt, whereas v2 gives \numScreeningDemoSitesGaTwoOThreeVaOTwoZeroVTwoMinEV{} to \numScreeningDemoSitesGaTwoOThreeVaOTwoZeroVTwoMaxEV{}~eV and is wrong by more, consistently. For SnO$_2$ the spread is the wider for v2 (\numScreeningDemoSitesSnOTwoVaOOneZeroVTwoMinEV{} to \numScreeningDemoSitesSnOTwoVaOOneZeroVTwoMaxEV{}~eV against \numScreeningDemoSitesSnOTwoVaOOneZeroVOneMinEV{} to \numScreeningDemoSitesSnOTwoVaOOneZeroVOneMaxEV{}~eV for v1). For Al$_2$O$_3$ the per-resample predictions of both backbones span far less than the error (v2: \numScreeningDemoSitesAlTwoOThreeVaOOneZeroVTwoMinEV{} to \numScreeningDemoSitesAlTwoOThreeVaOOneZeroVTwoMaxEV{}~eV, DFT \numScreeningDemoSitesAlTwoOThreeVaOOneZeroDftEV{}~eV), so that error persists when the training hosts change and does not come from the draw of the training set. Hypothesis: the structure-only network resolves the host-level stability and the coordination of a site better than the fine differences between sites of similar environment; ten sites cannot test this.

\subsubsection{What the demonstration shows.}
For \numScreeningDemoScopeNScreenedSites{} sites of \numScreeningDemoScopeNHosts{} oxides in the domain of the models, the structure-only network reproduces which oxides cost more energy to reduce, as two groups, and the order of three of the four within-host pairs; its mean absolute error is \numScreeningDemoSummaryVTwoMaeEV{}~eV (v2) and \numScreeningDemoSummaryVOneMaeEV{}~eV (v1). The wrong order in ZrO$_2$, the small predicted separation of the three-fold sites of Ga$_2$O$_3$ and the width of the intervals show that these data neither establish nor exclude a ranking of sites by energy finer than the error. A site of SnO$_2$ is shown as a structure render in the Supplementary Materials (\texttt{figures/si/structure\_sno2\_vacancy.png}): the oxygen of rutile-type SnO$_2$ has three cation neighbours in a plane, the coordination that Table~\ref{tab:screening_sites} records.
